\chapter{弗里德曼膨胀、辐射与暗物质}

\section{基本宇宙学方程回顾}

在比星系团更大的尺度上，宇宙被视为均匀的。选择坐标使我们位于\(x=0\)，并将一个代表性星系，或仅仅是一个随哈勃流运动的假想粒子，置于\(x=1\)。那么它到我们的物理距离按定义就是标度因子\(a(t)\)。位于\(x=2\)的点距离为\(2a(t)\)，依此类推。

\begin{figure}[ht]
\centering
\includegraphics[width=0.78\textwidth]{lecture_03_figure_02.png}
\caption{以\(x=0\)为中心的均匀球体、位于\(x=1\)的共动星系，以及\(a(t)\)的牛顿能量平衡。}
\end{figure}

只有半径\(a(t)\)的球体内部的物质才对位于\(x=1\)的粒子产生引力。给该粒子取单位质量。其动能为
\begin{equation}
\frac{1}{2}\dot a^2.
\end{equation}
球体所包围的质量为
\begin{equation}
M_{\mathrm{enc}}=\frac{4\pi}{3}a^3\rho,
\end{equation}
因此引力势能为
\begin{equation}
-\frac{G M_{\mathrm{enc}}}{a}
=
-\frac{4\pi G}{3}a^2\rho.
\end{equation}
\editorialnote{在讲座时间戳00:05:30处，黑板上清楚地显示了动能项和依赖于密度的势能项。此处写出的完整系数来自口头推导。}

能量守恒给出
\begin{equation}
\frac{1}{2}\dot a^2-\frac{4\pi G}{3}a^2\rho=\text{const}.
\end{equation}
该常数随后写为\(-k\)，并附有通常的提醒：标准归一化实际上只有在将方程乘以\(2\)之后才真正固定了该记号。在标准形式中，
\begin{equation}
\dot a^2=\frac{8\pi G}{3}a^2\rho-k,
\end{equation}
除以\(a^2\)得到
\begin{equation}
\left(\frac{\dot a}{a}\right)^2=\frac{8\pi G}{3}\rho-\frac{k}{a^2}.
\end{equation}
这就是整个讨论中所使用的弗里德曼方程形式。

\begin{figure}[ht]
\centering
\includegraphics[width=0.56\textwidth]{lecture_03_figure_03.png}
\caption{弗里德曼方程中分离出的曲率贡献\(-k/a^2\)。}
\end{figure}

曲率项在经验上非常接近于零，近到在这个层面上简单地忽略它是一个非常好的近似。它确实对应于空间曲率，但在此处从观测上看它太小，甚至无法确定其符号。

同一个方程也可以从广义相对论得到。以示意形式，时间-时间爱因斯坦方程为
\begin{equation}
R_{00}-\frac{1}{2}g_{00}R \propto T_{00}.
\end{equation}
\editorialnote{此处讨论只给出张量结构，而非精确的总体数值系数。}
对于膨胀的均匀度规，几何一侧产生\(\dot a/a\)与\(k/a^2\)的组合，而物质一侧产生能量密度。在这个阶段，牛顿论证和爱因斯坦方程说的是同一件事。

这也改变了\(\rho\)的含义。从现在起，最好将\(\rho\)称为能量密度。对于共动参考系中缓慢运动的物质，那只是静止能量密度。

\section{物质主导}

现在假设能量密度由共动参考系中静止的物质携带。取一个坐标边长为1单位的立方体。其物理边长为\(a(t)\)，因此其物理体积为\(a^3(t)\)。如果该共动立方体内的质能为\(M\)，且物质随流运动而不穿过边界，那么
\begin{equation}
\rho_m=\frac{M}{a^3}.
\end{equation}

这就是物质主导的情况。在这种理想化中，粒子在盒子中静止，因此温度取为零，压强取为零。压强需要粒子穿过盒子的边界。

取\(k\approx 0\)，方程变为
\begin{equation}
\left(\frac{\dot a}{a}\right)^2=\frac{8\pi G}{3}\frac{M}{a^3}.
\end{equation}
尝试幂律，
\begin{equation}
a(t)=Ct^p.
\end{equation}
那么
\begin{equation}
\frac{\dot a}{a}=\frac{p}{t},
\end{equation}
因此左侧按\(t^{-2}\)标度，而右侧按\(t^{-3p}\)标度。匹配幂次得到
\begin{equation}
3p=2,
\qquad
p=\frac{2}{3}.
\end{equation}
因此
\begin{equation}
a(t)\propto t^{2/3}.
\end{equation}
保留系数得到
\begin{equation}
C^3=6\pi G M.
\end{equation}

这个解的要点在于幂律。物质主导的宇宙按时间的三分之二次方膨胀，而非线性膨胀。该定律描述了一个漫长的中间时期，从大爆炸后大约\(10^5\)年到此后几十亿年。现今宇宙并非字面意义上的物质主导，因为暗能量现在占主导，但普通物质和宇宙历史的一大段时期正是以此方式描述的。

\section{辐射主导}

比那更早时，能量密度由辐射主导。原因很简单。宇宙包含数量巨大的光子，大约每个质子对应\(10^9\)个光子，而当宇宙被压缩时，这些光子变得更有能量。

光子没有静止质量，但它确实有能量。其能量为
\begin{equation}
E_\gamma=h\nu=\frac{hc}{\lambda}.
\end{equation}
关键的假设是，随着宇宙膨胀，光子波长也随之拉伸：
\begin{equation}
\lambda\propto a.
\end{equation}
盒子中的波给出了正确的图像。如果盒子增大，模式的波长随盒子一起增大。这里对膨胀宇宙使用了同样的想法。

由于 \(\lambda\propto a\)，每个光子的能量按如下方式标度：
\begin{equation}
E_\gamma\propto a^{-1}.
\end{equation}
再次取一个共动单位盒子。光子会进出，但平均而言，盒内的光子数保持不变。于是盒内光子的总能量也按 \(a^{-1}\) 标度。除以物理体积 \(a^3\)，辐射能量密度按如下方式标度：
\begin{equation}
\rho_r\propto a^{-4}.
\end{equation}

那个额外的 \(a^{-1}\) 因子就是全部差别所在。物质因为宇宙膨胀而被稀释。辐射也因为这个原因被稀释，同时还因为每个光子都被红移。

因此，辐射主导宇宙的弗里德曼方程为
\begin{equation}
\left(\frac{\dot a}{a}\right)^2=\frac{8\pi G}{3}\frac{M_r}{a^4}.
\end{equation}
\editorialnote{口头讨论中重复使用了符号 \(M\) 来表示这个新常数。这里加上下标 \(r\) 只是为了避免与物质情形混淆。}

使用同样的试探形式 \(a(t)=Ct^p\)。左边再次按 \(t^{-2}\) 标度，而右边按 \(t^{-4p}\) 标度。匹配幂次得到
\begin{equation}
4p=2,
\qquad
p=\frac{1}{2}.
\end{equation}
所以
\begin{equation}
a(t)\propto t^{1/2}.
\end{equation}
保留系数得到
\begin{equation}
C^4=\frac{32\pi G}{3}M_r.
\end{equation}

\begin{classroomqa}
\lecturetimestamp{00:35:41}
\audiencequestion{为什么那里会有一个 \(G\)？}
\lecturerresponse{因为产生引力作用的是能量密度，而不只是静质量。对于缓慢运动的物质，这一区别很容易被忽略，但光子迫使我们正视它：它们没有静质量，却仍然产生引力作用。}
\end{classroomqa}

一个装满辐射的盒子可以把这个论点具体化。推一个充满光子的盒子。一旦盒子加速，迎着光子运动的壁受到的撞击比背离光子运动的壁更强，因此必须更用力地推以对抗辐射压。于是盒子具有额外的惯性，而根据等效原理，这种额外惯性也意味着额外的引力质量。这就是这里用来论证能量产生引力作用的光子盒论证。

\begin{classroomqa}
\lecturetimestamp{00:38:11}
\audiencequestion{这两个方程中是否有任何一个可以用随时间减小的 \(a\) 求解？}
\lecturerresponse{有。如果恢复曲率项 \(k/a^2\)，就存在先膨胀后收缩的解。这并不是我们的宇宙看起来正在做的事，但方程允许这种解。}
\end{classroomqa}

\begin{classroomqa}
\lecturetimestamp{00:38:44}
\audiencequestion{能量守恒怎么了？}
\lecturerresponse{如果恢复 \(k\) 项，并把方程重新乘以 \(a^2\)，宇宙学方程就正是守恒定律。一侧的物质或辐射只是记账的一部分。膨胀贡献了动能部分和引力部分，而总平衡是守恒的。}
\end{classroomqa}

\section{物质与辐射并存}

真实的宇宙学同时包含物质和辐射。它们的密度按不同方式标度：
\begin{equation}
\rho_m\propto a^{-3},
\qquad
\rho_r\propto a^{-4}.
\end{equation}
所以在小的 \(a\) 下，辐射主导。在大的 \(a\) 下，物质主导。向前回溯，早期宇宙因此必定是辐射主导的。

\begin{figure}[ht]
\centering
\includegraphics[width=0.68\textwidth]{lecture_03_figure_05.png}
\caption{用于比较物质和辐射能量密度作为标度因子函数的坐标轴。}
\end{figure}

\editorialnote{在讲座时间戳 00：44：45 处，黑板上只显示了比较图起始处的坐标轴。紧接着讨论的物质和辐射曲线在正文中描述，这里不再重绘。}

物质曲线按 \(a^{-3}\) 下降。辐射曲线按 \(a^{-4}\) 下降。两者在 \(a\to 0\) 时都发散，但辐射曲线在那里上升得更快，所以两条曲线在某个相等标度因子处相交。粗略估计把这一转变放在大约 \(10^4\) 年之后，也许稍多一点。

同一点也可以用 \(a(t)\) 对 \(t\) 的图来表示。早期，膨胀表现得像 \(t^{1/2}\)。后来它表现得像 \(t^{2/3}\)。这种变化不是跳跃；它是通过一个相对狭窄的时间范围发生的连续过渡。

\begin{classroomqa}
\lecturetimestamp{00:44:04}
\audiencequestion{随着宇宙变老，幂 \(p\) 会从 \(1/2\) 连续地变为 \(2/3\) 吗？}
\lecturerresponse{是的。它不会跳跃。这种变化是连续的，尽管它发生在一个相当狭窄的时间范围内。}
\end{classroomqa}

\section{温度、复合与最后散射}

下一个量是温度。如果一个典型光子的能量按 \(a^{-1}\) 标度，那么辐射浴的温度也如此：
\begin{equation}
T\propto a^{-1}.
\end{equation}
如果 \(a\) 加倍，温度就减半。

取当前光子温度约为 \(3\,\mathrm K\)，电离温度约为 \(3000\,\mathrm K\)，大致是太阳表面的温度。那么
\begin{equation}
\frac{a_{\mathrm{today}}}{a_{\mathrm{ion}}}\sim 10^3.
\end{equation}
使用晚期物质主导定律 \(a\propto t^{2/3}\)，
\begin{equation}
\left(\frac{t_{\mathrm{today}}}{t_{\mathrm{ion}}}\right)^{2/3}\sim 10^3,
\end{equation}
所以
\begin{equation}
\frac{t_{\mathrm{today}}}{t_{\mathrm{ion}}}\sim 1000^{3/2}\sim 3\times 10^4.
\end{equation}
在这个估计的精度下，取宇宙年龄为 \(10^{10}\) 年，得到
\begin{equation}
t_{\mathrm{ion}}\sim 3\times 10^5\ \mathrm{years}.
\end{equation}

这是氢和氦不再完全电离的时期。在此之前，宇宙是炽热的等离子体。为什么等离子体是不透明的？因为自由电荷对电磁波有强烈响应。电场驱动电荷来回运动，介质表现得像导体，而导体是不透明的。这就是与太阳表面相关的类比。

一旦宇宙冷却到足以让电子与原子核复合，普通中性原子便形成，宇宙变得透明得多。这一转变并非无限锐利。它是一个模糊的位置，或者说一个模糊的时间区间，但对于宇宙学而言已经足够锐利，因此被称为最后散射面。

向外看意味着向过去看。邻近星系展示的是最近的过去。更遥远的星系展示的是更早的时期。足够远的地方，每条视线都到达宇宙仍然电离且不透明的时期。电磁辐射无法展示比那更早的任何东西。

其他信使原则上可以探测更早的时期。中微子电中性，会更早退耦。引力子，如果人们能够观测到的话，将来自更早的时期。但是，让它们能够穿过早期等离子体的那种相互作用的微弱性，也使它们极难被探测到。

如果宇宙曾经在每个方向上都像太阳表面一样炽热，为什么天空不是令人目眩地明亮？因为最后散射面正在退行，而从它到达的光子已经被红移了大致相同的因子\(10^3\)，这个因子把\(3000\,\mathrm K\)与\(3\,\mathrm K\)分开。到达我们的是微波辐射，而不是可见的阳光。

\begin{classroomqa}
\lecturetimestamp{01:03:51}
\audiencequestion{所以我们的望远镜无法看到大约\(300{,}000\)年之前？}
\lecturerresponse{电磁望远镜不能。在那之前，宇宙是电离且不透明的。其他信使，至少原则上，可以探测更早的时期。}
\end{classroomqa}

\begin{classroomqa}
\lecturetimestamp{01:08:43}
\audiencequestion{当你谈论光子的温度时，你指的是具有那个峰值辐射的黑体的温度。但那不是与波长线性相关的，对吧？}
\lecturerresponse{对。温度本质上是典型光子的能量，所以对于光子来说，它与波长成反比。\(T^4\)定律指的是能量密度，而不是单个光子的能量。它把每个光子的能量与单位体积内的光子数结合起来，而数密度按\(T^3\)缩放。}
\end{classroomqa}

因此
\begin{equation}
n_\gamma\propto T^3,
\qquad
\rho_\gamma\propto T^4.
\end{equation}

\section{来自旋转曲线的暗物质}

暗物质是通过星系旋转曲线进入的。它不是从大尺度弗里德曼方程推断出来的，而是从星系尺度上的天文学推断出来的。

从最简单的牛顿问题开始：一个质量为\(m\)的轻物体在半径\(M\)处绕中心质量\(r\)运行。引力为
\begin{equation}
\frac{GMm}{r^2},
\end{equation}
而向心加速度为\(v^2/r\)。牛顿定律给出
\begin{equation}
m\frac{v^2}{r}=\frac{GMm}{r^2}.
\end{equation}
消去\(m\)，乘以\(r\)，得到
\begin{equation}
v(r)=\sqrt{\frac{GM}{r}}.
\end{equation}
因此，如果引力质量集中在中心附近，轨道速度就像\(r^{-1/2}\)那样下降。

如果质量跟随光，这就是人们所预期的。星系中大多数可见物质位于核心附近，因此外围恒星的行为会很像外太阳系中的行星。天文学家通过测量邻近星系两侧的多普勒频移差异来推断轨道速度，在减去星系整体的退行运动后，比较朝向我们运动的一侧与远离我们运动的一侧。测量到的轨道速度在大半径处大致保持恒定。

因此把中心质量\(M\)替换为包含质量\(M(r)\)：
\begin{equation}
\frac{v^2}{r}=\frac{G M(r)}{r^2}.
\end{equation}
如果\(v(r)\)近似恒定，那么
\begin{equation}
M(r)\propto r.
\end{equation}
换句话说，即使在可见光不增加的地方，引力质量也随着半径持续增加。

现在假设，正如讨论为简单起见所做的那样，质量分布大致是球形的。那么
\begin{equation}
M(R)=\int_0^R 4\pi r^2\rho(r)\,dr.
\end{equation}
如果\(M(R)\propto R\)，对\(R\)求导给出
\begin{equation}
4\pi R^2\rho(R)\propto 1,
\end{equation}
因此
\begin{equation}
\rho(r)\propto \frac{1}{r^2}.
\end{equation}

\begin{classroomqa}
\lecturetimestamp{01:28:18}
\audiencequestion{看起来密度应该像\(1/r^2\)那样。这有什么原因吗？}
\lecturerresponse{从观测到的平坦旋转曲线可以推断出\(M(R)\propto R\)，求导后给出\(\rho(r)\propto 1/r^2\)。除此之外，这里不声称有简单的第一性原理解释；详细的模拟完成了大部分工作。}
\end{classroomqa}

稳定的结论是，一个星系所含的引力质量远多于仅由发光物质所能推得的质量，大约要大数倍，而按这里所用的估计则大致大一个数量级。可见星系位于一个远为巨大的暗晕之中。

是什么使普通气体无法停留在如此弥散的状态？是耗散。气体云之所以收缩，是因为碰撞、类摩擦过程以及辐射使它得以损失能量。暗物质显然并不显著辐射，而且几乎是无碰撞的。引力可以把它向内拉，并将势能转化为动能，但若没有有效的耗散，它就不会坍缩成普通物质所形成的那种致密发光结构。使暗晕保持延展的是动能。

随后谨慎地讨论了几种替代方案。据说中微子速度太快，无法在星系尺度上如此表现。由普通氢构成的暗晕会严重破坏早期宇宙核合成。由普通物质构成的黑洞也被提到是有问题的，因为它们会意味着早期宇宙中有过多的普通物质，并会引发引力透镜方面的问题；讨论在此处并未声称超出这一点的结论。

暗晕并未被描述为薄壳或环面。它是一个近似球形、非均匀的球体，越靠近中心越致密，向外则逐渐变稀。

\begin{classroomqa}
\lecturetimestamp{01:53:12}
\audiencequestion{你是在说子弹星系团吗？}
\lecturerresponse{是的。X射线气体显示的是普通物质，而引力透镜显示的则是总引力质量。在碰撞中，这两种分布分开了。这是强有力的证据，表明存在真实的暗物质，而不是仅仅与可见物质相关联的单纯引力修改。}
\end{classroomqa}

这就是最后的观测要点。关于暗物质究竟是什么，人们对其细节所知甚少，但这里没有任何内容需要一条激进的新原理。保守的可能性仅仅是：自然界包含尚未被识别出的额外弱相互作用粒子。
